\pagebreak
\sectiontitle{Beschaffung}

\formheading{ABC-Analyse} \; $\text{Anteil}_i = \dfrac{\text{Volumen}_i}{\text{Gesamtvolumen}} \cdot 100\,\%$, \quad $\text{Volumen}_i = \text{Menge}_i \times \text{Preis}_i$

\formheading{Scoring-Methode / Nutzwertanalyse} \; $\text{Gesamtnutzwert} = \sum (\text{Gewicht} \times \text{Bewertung})$

\formheading{Kostenvergleichsrechnung (KVR)}
\begin{minipage}[c]{0.56\textwidth}
\centering
\[ K = B + \dfrac{A}{t} + \dfrac{A}{2} i \;\text{(o. Restwert)}, \quad K = B + \dfrac{A - R}{t} + i \dfrac{A + R}{2} \;\text{(m. Restwert)}, \quad k = \dfrac{K}{x} \]
\end{minipage}%
\hfill
\begin{minipage}[c]{0.38\textwidth}
\centering
\begin{symtable}[3cm]
	$B,A,R$ & Betriebs-, Anschaffungskosten, Restwert \\
	$t,i,x$ & Nutzungsdauer, Kalk.-zins, Mengenoutput \\
\end{symtable}
\end{minipage}

\formheading{Optimale Bestellmenge -- Andler-Formel (EOQ)}
\begin{minipage}[c]{0.56\textwidth}
\centering
\[ h = \alpha c, \quad K_{\mathrm{L}} = \dfrac{q}{2} h, \quad K_{\mathrm{B}} = \dfrac{r}{q} s, \quad K_{\mathrm{ges}} = \dfrac{q}{2} h + \dfrac{r}{q} s \]
\[ q^{*} = \sqrt{\dfrac{2 r s}{h}} \;\text{(Optimum)}, \quad T = \dfrac{q}{r},\ T^{*} = \dfrac{q^{*}}{r} \;\text{(Reichw.)}, \quad U = c \cdot r \]
\end{minipage}%
\hfill
\begin{minipage}[c]{0.38\textwidth}
\centering
\begin{symtable}[3cm]
	$r,c$ & Bedarf [ME/J], Einkaufspreis [EUR/ME] \\
	$h,s$ & Lagerkosten $=\alpha c$, Bestellkosten \\
	$q,q^{*},\alpha$ & (opt.) Bestellmenge, Lagerzins \\
\end{symtable}
\end{minipage}

\formheading{Lagerbestandsformeln}
\[ \begin{aligned}
\varnothing\,\text{Arbeitsbestand} &= \frac{q}{2} \\
\varnothing\,\text{Gesamtbestand} &= \varnothing\,\text{Arbeitsbestand} + \varnothing\,\text{Sicherheitsbestand} \\
\varnothing\,\text{Bestand} &= \varnothing\,\text{Flussstärke} \cdot \Delta t \\
\text{Buchbestand} &= \text{Inventurbestand} + \text{Zugänge} - \text{Abgänge} \\
\text{Verfügb. Bestand} &= \text{phys. Bestand} + \text{Bestellbestand} - \text{reserv. Bestand}
\end{aligned} \]
